% Scala, JavaScript, Python, Ruby

\subsection{Metrics Evaluation for Scala}

Scala is a hybrid object-functional language with a strict type system and advanced object-oriented mechanisms.
It aims to combine object-oriented and functional programming through a uniform object model and extended types~\cite{odersky2004overview, emir2007matching}.
In this section, Scala is evaluated according to the proposed metrics based on the language specification and scientific papers.

\subsubsection{Type System and Subtyping Model}

\textbf{1) Nominal–Structural Typing Score $S_{nom}(L)$}

Scala employs nominal subtyping through explicit declarations such as \texttt{class}, \texttt{trait}, and inheritance clauses (e.g., \texttt{extends}), where subtype relations are defined by named entities and governed by formal typing judgments. 
Scala also supports structural types via refinement types, where type compatibility is determined by the presence of members rather than explicit inheritance, and this behaviour is formally specified in its typing rules \cite{odersky2004overview}. 

\[
    S_{nom}(\text{Scala}) = 1, I_{struct} = 1 \Rightarrow S_{nom} = 1
\]

\textbf{2) Variance Formalization Score $S_{var}(L)$}

Scala provides fully explicit variance annotations, where type parameters can be declared covariant (\texttt{+A}), contravariant (\texttt{-A}), or invariant (default), and their usage is syntactically enforced by the compiler through well-defined typing rules.

$V_{exp}(\text{Scala}) = 1$. 

Furthermore, the specification formalizes variance soundness via restrictions on positions of type parameters (e.g., covariant types cannot appear in contravariant positions), ensuring type safety through compile-time checks derived from formal subtyping and variance rules.

$V_{sound}(\text{Scala}) = 1$.

\[
    S_{var}(\text{Scala}) = 1
\]

\textbf{3) Inheritance–Subtyping Separation Score $S_{rel}(L)$}

First, although nominal subtyping is primarily introduced via class and trait inheritance, the language also permits limited explicit subtyping without direct inheritance through constructs such as refinement types and type bounds (e.g., $A <: B$ in type parameters), indicating partial but not fully general separation.

$I_{dec}(\text{Scala})=0.5$

Second, Scala’s subtyping relation is formally richer than inheritance alone (including variance, compound types, and path-dependent types), yet still partially grounded in declared inheritance hierarchies, implying only partial independence.

$I_{ind}(\text{Scala})=0.5$

Third, the specification distinguishes subtyping ($<:$) from inheritance (class/trait extension) at a conceptual level, but does not fully axiomatize them as completely orthogonal relations. 

$I_{form}(\text{Scala})=0.5$

$S_{rel}(\text{Scala})=\frac{0.5+0.5+0.5}{3}=0.5$

\textbf{Final metric value:}

\[
\begin{aligned}
    TSSI(\text{Scala}) = 0.3 \cdot S_{nom}(\text{Scala}) + 0.3 \cdot S_{var}(\text{Scala}) + 0.4 \cdot S_{rel}(\text{Scala})\\
    = 0.3 \cdot 1 + 0.3 \cdot 1 + 0.4 \cdot 0.5 \\
    = 0.8
\end{aligned}
\]

\subsubsection{Abstraction and Encapsulation Mechanisms}

\textbf{1) Visibility Granularity Index $V(L)$}

Scala provides multiple semantically distinct access control levels: \texttt{private[this]} (object-private), \texttt{private} (class-private), qualified \texttt{private[pkg]} (package-restricted), \texttt{protected[this]} (object-protected), \texttt{protected} (subclass access), and unrestricted public access (default).

\[
    n_{levels}(\text{C\#})=7 \Rightarrow V(\text{Scala}) = \frac{6}{7}
\]


\textbf{2) Module Isolation Strength $M(L)$}

Scala has 5 ways to organize code: packages, package objects, singleton objects, classes, and traits.

$b_{total}(\text{Scala}) = 5$

The language does not require you to explicitly specify which elements are exported from the module. Instead, access is regulated by visibility modifiers (for example, `private[pkg]`). Therefore, there are no strict modular boundaries in Scala with a mandatory explicit export interface.

$b_{strict}(\text{Scala}) = 0$

\[
    M(\text{Scala}) = \frac{0}{5} = 0
\]


\textbf{3) Interface Abstraction Capacity $I(L)$}

Scala separates specification and implementation using two main mechanisms: `trait` and `abstract class'. The `trait` is used for abstractions and mixin composition, while the `abstract class` is used for partial implementation with abstract elements.

$n_{abstraction\_constructs}(\text{Scala}) = 2$

The set of object-defining constructs includes \texttt{class}, \texttt{trait}, and \texttt{object} (singleton).

$n_{object}(\text{Scala}) = 3$

\[
    I(\text{Scala}) = \frac{2}{3} \approx 0.67
\]


\textbf{4) Encapsulation Enforcement Coefficient $E(L)$}

Scala enforces encapsulation at the type system level via access modifiers (e.g., \texttt{private}, \texttt{protected}, and qualified visibility), it does not fully prevent encapsulation violations at runtime. Specifically, Scala inherits unrestricted reflection capabilities from the Java ecosystem (e.g., \texttt{java.lang.reflect}) and allows access to private members through reflective APIs, as well as other unsafe mechanisms available on the JVM. These capabilities are not restricted by the Scala type system and are not prohibited in the specification. 

\[
    E(\text{Scala}) = 0
\]

\textbf{Final metric value:}

\[
\begin{aligned}
    AE(\text{Scala}) = 0.2 \cdot V(\text{Scala}) + 0.25 \cdot M(\text{Scala}) + 0.25 \cdot I(\text{Scala}) + 0.3 \cdot E(\text{Scala}) \\
    = 0.2 \cdot \frac{6}{7} + 0.25 \cdot 0 + 0.25 \cdot \frac{2}{3} + 0.3 \cdot 0
    \approx 0.3381
\end{aligned}
\]

\subsubsection{Inheritance Semantics and Hierarchy Management}

\textbf{1) Hierarchy Simplicity Factor $H_s(L)$}

Scala enforces single inheritance for classes (each class may extend at most one superclass) while allowing multiple inheritance only through \texttt{trait} mixins, which are composed using a well-defined linearization algorithm that preserves a single inheritance chain at the class level. Thus, the maximum number of direct superclass declarations for classes is 1.

\[
    H_s(\text{Scala}) = 1
\]


\textbf{2) Linearization Determinism Score $L_d(L)$}

Scala uses a strict and predictable algorithm for determining the method invocation order for classes and traits.
It forms a single inheritance hierarchy, which ensures unambiguous resolution of redefined methods and eliminates ambiguities in multiple inheritance \cite{odersky2004overview}.

\[
    L_d(\text{Scala}) = 1
\]

\textbf{3) Constraint Protection Coefficient $C_p(L)$}

$p_1(\text{Scala})=1$ since the language provides \texttt{final} for classes and methods, preventing further inheritance or overriding

$p_2(\text{Scala})=1$ due to \texttt{sealed} traits/classes restricting subclassing to the same compilation unit

$p_3(\text{Scala})=1$ because overriding requires the explicit \texttt{override} modifier, enforced at compile time

$p_4(\text{Scala})=0$ as methods are virtual by default and dynamic dispatch is the standard semantics

$p_5(\text{Scala})=0$ because the language does not enforce a strict separation between abstract classes and interfaces, instead unifying them via traits

$p_6(\text{Scala})=1$ because the compiler enforces inheritance constraints such as prohibiting extension of \texttt{final} classes and enforcing linearization rules

\[
    C_p(\text{Scala}) = \frac{4}{6} \approx 0.67
\]

\textbf{4) Fragile Base Mitigation Factor $F_b(L)$}

Scala incorporates several semantic safeguards that mitigate fragile base class effects. These include mandatory \texttt{override} annotations for redefining inherited members, fine-grained visibility control (e.g., \texttt{private}, \texttt{protected}, and qualified access), restrictions on variance positions ensuring type safety, and a well-defined linearization model governing \texttt{super} calls in mixin compositions.

$s_1(\text{Scala})=1$ \texttt{override} modifier is mandatory and prevents accidental overriding

$s_2(\text{Scala})=1$ the language provides strict visibility modifiers (e.g., \texttt{private}, \texttt{protected}, and scoped qualifiers) that control exposure of overridable members.

$s_3(\text{Scala})=0$ explicit superclass invocation is not universally required and overriding methods may omit \texttt{super} calls.

$s_4(\text{Scala})=1$ variation annotations ($+$/$-$) provide type-safe subtyping and overriding.

$s_5(\text{Scala})=0$ default methods can be overridden if they are not declared \texttt{final}.

$s_6(\text{Scala})=1$ the compiler enforces strict override consistency checks on method signatures and types.

\[
    F_b(\text{Scala})=\frac{1+1+0+1+0+1}{6}=\frac{4}{6}\approx 0.67
\]

\textbf{Final metric value:}
\[
\begin{aligned}
    ISRI(\text{Scala}) = 0.2 \cdot H_s(\text{Scala}) + 0.25 \cdot L_d(\text{Scala}) + 0.3 \cdot C_p(\text{Scala}) + 0.25 \cdot F_b(\text{Scala}) \\
    = 0.2 \cdot 1 + 0.25 \cdot 1 + 0.3 \cdot \frac{4}{6} + 0.25 \cdot \frac{4}{6}
    = 0.8167
\end{aligned}
\]

\subsubsection{Composition and Alternatives to Inheritance}

\textbf{1) Class-Based Delegation Score $D_{cls}(L)$}

Scala does not have a built-in delegation mechanism or a special keyword. Delegation is usually implemented through composition and manual forwarding of method calls. There are also no automatic methods forwarding at the language level.

\[
    K_{del}(\text{Scala}) = 0.5, S_{fwd}(\text{Scala}) = 0 \Rightarrow D_{cls}(\text{Scala}) = 0.25
\]

\textbf{2) Trait Expressiveness Score $T_{exp}(L)$}

Scala resolves name conflicts during trait composition using a deterministic linearization order. In cases of ambiguity, the rightmost trait in the mixin order takes precedence, and conflicts are resolved through overriding rules defined in the specification. 
Furthermore, Scala traits can declare requirements on the classes that mix them in via self-types (e.g., \texttt{this: SomeType =>}), which are formally specified constraints on the composition context.

\[
    R_{conf}(\text{Scala}) = 0.5, O_{mod}(\text{Scala}) = 1 \Rightarrow R_{conf}(\text{Scala}) = 0.75
\]

\textbf{3) Interface Composition Capacity $I_{comp}(L)$}

Scala does not have a separate \texttt{interface} construct. Instead, \texttt{trait}s serve as their direct analogue and extend beyond traditional interfaces. A class may mix in an unrestricted number of traits using the \texttt{with} keyword. 
Traits can contain fully defined method implementations, fields, and even initialisation logic, as formally specified in the language.

\[
    N_{impl}(\text{Scala}) = 1, D_{def}(\text{Scala}) = 1 \Rightarrow I_{comp} = 1
\]

\textbf{Final metric value:}

\[
\begin{aligned}
    CRFI(\text{Scala}) = 0.3 \cdot D_{cls}(\text{Scala}) + 0.35 \cdot T_{exp}(\text{Scala}) + 0.35 \cdot I_{comp}(\text{Scala}) \\
    = 0.3 \cdot 0.25 + 0.35 \cdot 0.75 + 0.35 \cdot 1 = 0.6875
\end{aligned}
\]

\subsubsection{Object Representation and Creation Model}

\textbf{1) Identity Semantics Score $S_{id}(L)$}

All user-defined types (classes, traits) are reference types:
Scala distinguishes between reference types (instances of \texttt{class}, \texttt{trait}, and \texttt{object}) and value types (e.g., subclasses of \texttt{AnyVal}). But all user-definable type are reference types

\[
    S_{id} = 1
\]

\textbf{2) Instantiation Paradigm Diversity $S_{cr}(L)$}

Scala natively supports standard constructor-based instantiation via the \texttt{new} keyword, singleton object instantiation through \texttt{object} declarations (module pattern with implicit initialisation), and literal construction for built-in and case classes (e.g., \texttt{apply} methods enabling factory-style creation without \texttt{new}). Additionally, case classes provide a canonical factory mechanism via automatically generated companion objects.

\[
    P_{supported}(\text{C\#}) = 4 \Rightarrow S_{cr}(\text{Scala}) = \frac{3}{4}
\]

\textbf{3) Meta-level Extensibility Score $S_{meta}(L)$}

Scala supports runtime introspection through reflection APIs (e.g., \texttt{scala.reflect} and interoperability with \texttt{java.lang.reflect}), allowing programs to query type information and object structure dynamically. However, the language does not natively support full structural modification of objects (e.g., adding/removing fields or methods at runtime) nor overriding the object instantiation protocol via metaobjects within the core specification.

\[

    I_{level} = 1, I_{level}(\text{JS/Python/Ruby}) = 3 \Rightarrow S_{meta}(\text{Scala}) = \frac{1}{3}
\]

\textbf{Final metric value:}
\[
\begin{aligned}
    ORCI(\text{Scala}) = 0.35 \cdot S_{id}(\text{Scala}) + 0.3 \cdot S_{cr}(\text{Scala}) + 0.35 \cdot S_{meta}(\text{Scala}) \\
    = 0.35 \cdot 1 + 0.3 \cdot \frac{3}{4} + 0.35 \cdot \frac{1}{3} \approx 0.6917
\end{aligned}
\]

\subsubsection{Polymorphism and Dispatch Mechanisms}

\textbf{1) Dispatch Scope Mechanism $M_{scope}(L)$}

Scala natively supports single dispatch based on the dynamic type of the receiver object.

$I_{single}(\text{Scala}) = 1$

But, the language does not provide true multiple dispatch in the sense of dynamic selection based on all argument types. Method overloading is resolved statically at compile time rather than dynamically.

$I_{multi}(\text{Scala}) = 0$

Scala does not support predicate-based dispatch as a built-in mechanism. The choice of run-time behaviour is usually implemented through pattern matching or flow control constructs.

$I_{pred}(\text{Scala}) = 0$

\[
    M_{scope}(\text{Scala}) = \frac{1}{5} = 0.2
\]

\textbf{2) Dispatch Safety Guarantee $S_{disp}(L)$}

Scala statically verifies the correctness of method invocation through the type system: method resolution, redefinition correctness, and secure dynamic dispatch.

$I_{static}(\text{Scala})=1$

Scala verifies the completeness of pattern matching for \texttt{sealed} hierarchies and algebraic data types. The compiler warns or reports errors if not all cases are handled.

$I_{exhaust}(\text{Scala})=1$

\[
    S_{disp}(\text{Scala}) = \frac{1+1}{2} = 1
\]

\textbf{3) Virtualisation Optimization Potential $O_{virt}(L)$}

Scala supports \texttt{final} declarations for classes and methods, preventing overriding and enabling deVirtualisation opportunities.

$K_{final}(\text{Scala}) = 1$

It also provides \texttt{sealed} classes and traits, which restrict subclassing to the same compilation unit and allow the compiler to reason about closed hierarchies.

$K_{sealed}(\text{Scala}) = 1$

Scala does not define explicit static methods or language-level static dispatch annotations in the same sense as some other languages (despite JVM-level optimizations), and dispatch remains primarily virtual unless restricted by \texttt{final}.

$K_{static}(\text{Scala}) = 0$

\[
    O_{virt}(\text{Scala}) = \frac{1 + 1 + 0}{3} \approx 0.67
\]

\textbf{Final metric value:}
\[
\begin{aligned}
    PDEI(\text{Scala}) = 0.35 \cdot M_{scope}(\text{Scala}) + 0.35 \cdot S_{disp}(\text{Scala}) + 0.3 \cdot O_{virt}(\text{Scala}) \\
    = 0.35 \cdot 0.2 + 0.35 \cdot 1 + 0.3 \cdot \frac{2}{3} = 0.62
\end{aligned}
\]

\subsubsection{State Model, Mutability, and Concurrency Guarantees}

\textbf{1) Default Mutability Score $M_{def}(L)$}

although Scala promotes immutability idiomatically (e.g., via \texttt{val} bindings and immutable collections), mutability is explicitly available and controlled through the \texttt{var} keyword and mutable collection types, meaning that immutability must be declared rather than being enforced by default at the object level.
\[
    M_{def}(\text{Scala}) = 0.5
\]

\textbf{2) Static Reference Constraint Density $S_{dens}(L)$}

First, the constructs that introduce or manipulate references in Scala programming language: object references (all values are references except primitives under erasure), \texttt{val} bindings, \texttt{var} bindings, class/trait fields, parameter passing (by-value references), and path-dependent and refined types that qualify references.

$N_{ref}(\text{Scala})=6$

Second, the constructs that enforce static restrictions: \texttt{val} (immutability of the reference binding), path-dependent types (restricting reference usage via stable paths), and refined/structural types (constraining accessible members at compile time), while \texttt{var}, general object references, fields, and parameter passing do not impose strict aliasing or mutability constraints.

$N_{constr}(\text{Scala})=3$

\[
    S_{dens}(\text{Scala}) = \frac{3}{6} = 0.5
\]

\textbf{3) Concurrency Primitive Safety $P_{conc}(L)$}

First, Scala’s type system does not enforce data race freedom (e.g., shared mutable state via \texttt{var} is permitted without compile-time guarantees), and the Java Memory Model underlies its execution semantics. The correctness relies largely on programmer discipline and external libraries rather than intrinsic type safety.

$S_{race}(\text{Scala}) = 0$

Second, while the core language exposes low-level JVM threads, Scala’s standard library provides \texttt{Future} and \texttt{Promise} abstractions (structured concurrency), and widely adopted frameworks such as Akka introduce Actor-based high-level concurrency models. However, since these are not strictly part of the core language specification.

$S_{model}(\text{Scala}) = 0.5$.

\[
    P_{conc}(\text{Scala}) = 0.25
\]

\textbf{Final metric value:}
\[
\begin{aligned}
    SCSI(\text{Scala}) = 0.35 \cdot M_{def}(\text{Scala}) + 0.25 \cdot S_{dens}(\text{Scala}) + 0.4 \cdot P_{conc}(\text{Scala}) \\
    = 0.35 \cdot 0.5 + 0.25 \cdot 0.5 + 0.4 \cdot 0.25 = 0.4
\end{aligned}
\]

\subsubsection{Support for Formal Specification and Verifiability}

\textbf{1) Native Contract Constructs Score $C_{nat}(L)$}

$n_{pre}(\text{Scala}) = 0$: there are no built-in prerequisites in Scala; \texttt{require} is a library method, not a language mechanism.

$n_{post}(\text{Scala}) = 0$: the language does not support built-in postconditions; \texttt{ensuring} is also implemented only as a library method.

$n_{inv}(\text{Scala}) = 0$: Scala does not have built-in support for class invariants.

\[
    C_{nat}(\text{Scala}) = 0
\]

\textbf{2) Verification Enforcement Level $V_{enf}(L)$}

Scala does not support built-in static verification or formal proofs in the compiler. However, the standard library provides runtime verification mechanisms such as \texttt{require}, \texttt{assert}, and \texttt{ensuring}.

\[
    V_{enf}(\text{Scala}) = 0.5
\]

\textbf{Final metric value:}

\[
\begin{aligned}
    FSVI(\text{Scala}) = 0.45 \cdot C_{nat}(\text{Scala}) + 0.55 \cdot V_{enf}(\text{Scala}) \\
    = 0.45 \cdot 0 + 0.55 \cdot 0.5 = 0.275
\end{aligned}
\]

\subsection{Metrics Evaluation for JavaScript}

JavaScript is a dynamically typed object-oriented language standardised by ECMAScript. Unlike languages with static typing and classes, it uses prototypal inheritance, dynamic property lookup, and late binding. Objects in JavaScript are sets of properties\cite{ecma262}, and the language is focused on flexibility and execution during program operation.

\subsubsection{Type System and Subtyping Model}

\textbf{1) Nominal–Structural Typing Score $S_{nom}(L)$}

JavaScript does not provide a formal static type system with structural subtyping rules. It employs dynamic typing where type compatibility is resolved at runtime without formal typing judgments or compile-time subtyping relations. Although object shapes can be compared implicitly (duck typing), this behaviour is not specified via formal structural typing rules in the language specification. 

JavaScript lacks nominal typing: there are no explicit subtype declarations, and the prototype-based inheritance model does not constitute nominal subtyping in the formal type-theoretic sense. Subtyping relationships are not enforced through named type declarations but rather through dynamic property lookup. 

\[
    I_{struct}(\text{JavaScript}) = 0, I_{nom}(\text{JavaScript}) = 0 \Rightarrow S_{nom}(\text{JavaScript}) = 0
\]

\textbf{2) Variance Formalization Score $S_{var}(L)$}

ECMAScript does not include syntactic constructs for declaring parametric polymorphism or variance annotations (such as covariance or contravariance), and all values are dynamically typed without compile-time type parameterization. 
There are no formally defined variance rules, nor any specification-level guarantees of type soundness related to variance. The type behaviour is governed entirely by runtime semantics without static verification or proofs of safety. 

\[
    V_{exp}(\text{JavaScript}) = 0, V_{sound}(\text{JavaScript}) = 0 \Rightarrow S_{var}(\text{JavaScript}) = 0
\]

\textbf{3) Inheritance–Subtyping Separation Score $S_{rel}(L)$}

First, JavaScript allows subtyping (i.e., object compatibility and behavioural substitutability) independently of explicit inheritance constructs: objects can be created via object literals or functions and still be used interchangeably based on their structure, without requiring the \texttt{[[Prototype]]} linkage.

$I_{dec}(\text{JavaScript})=1$

Second, subtyping is not formally defined in the specification as an independent relation, but emerges implicitly through structural behaviour (duck typing) and internal methods (e.g., \texttt{[[Get]]}, \texttt{[[Set]]}). The inheritance via the prototype chain influences property resolution, making the independence partial. 

$I_{ind}(\text{JavaScript})=0.5$

Third, the ECMAScript specification does not formally distinguish subtyping from inheritance: it defines inheritance operationally via prototype chains, while subtyping is not explicitly formalized as a separate semantic relation.

$I_{form}(\text{JavaScript})=0$

\[
    S_{rel}(\text{JavaScript}) = \frac{1 + 0.5 + 0}{3} = 0.5.
\]

\textbf{Final metric value:}
\[
\begin{aligned}
    TSSI(\text{JavaScript}) = 0.3 \cdot S_{nom}(\text{JavaScript}) + 0.3 \cdot S_{var}(\text{JavaScript}) + 0.4 \cdot S_{rel}(\text{JavaScript}) \\
    = 0.3 \cdot 0 + 0.3 \cdot 0 + 0.4 \cdot 0.5 = 0.2
\end{aligned}
\]

\subsubsection{Abstraction and Encapsulation Mechanisms}

\textbf{1) Visibility Granularity Index $V(L)$}

JavaScript supports: public visibility for all object properties by default, class-private fields (denoted with the \# syntax), which are enforced at the specification level via lexical scoping and inaccessible outside the declaring class\cite{ecma262}.

\[
    n_{levels}(\text{JavaScript}) = 2, n_{levels}(\text{C\#})=7 \Rightarrow
    V(\text{JavaScript}) = \frac{2}{7}
\]

\textbf{2) Module Isolation Strength $M(L)$}

JavaScript uses ES modules as the main modularity mechanism. Data is accessed between modules only through \texttt{export} and \texttt{import}, which provides isolation and access control. Scopes of scripts and functions are not considered a full-fledged modular system.

\[
    M(\text{JavaScript}) = 1
\]


\textbf{3) Interface Abstraction Capacity $I(L)$}

JavaScript does not have built-in abstraction tools such as interfaces or abstract classes. Abstraction is usually implemented through conventions, such as duck typing or documentation.

\[
    I(\text{JavaScript}) = 0
\]


\textbf{4) Encapsulation Enforcement Coefficient $E(L)$}

JavaScript allows you to dynamically access the properties of objects and change them at runtime. The \texttt{Reflect} and \texttt{Proxy} mechanisms make it possible to bypass encapsulation restrictions for all members.

\[
    E(\text{JavaScript}) = 0
\]

\textbf{Final metric value:}
\[
\begin{aligned}
    AE(\text{JavaScript}) = 0.2 \cdot V(\text{JavaScript}) + 0.25 \cdot M(\text{JavaScript}) + 0.25 \cdot I(\text{JavaScript}) + 0.3 \cdot E(\text{JavaScript}) \\
    = 0.2 \cdot \frac{2}{7} + 0.25 \cdot 1 + 0.25 \cdot 0 + 0.3 \cdot 0 \\
    \approx 0.3071
\end{aligned}
\]

\subsubsection{Inheritance Semantics and Hierarchy Management}

\textbf{1) Hierarchy Simplicity Factor $H_s(L)$}

JavaScript supports only single inheritance in terms of object delegation: each object has at most one internal \texttt{[[Prototype]]} reference, and class-based syntax (\texttt{extends}) likewise allows only one direct superclass. Although mixin patterns can simulate multiple inheritance, they are not formal language-level inheritance mechanisms defined in the specification.

\[
    H_s(\text{JavaScript}) = 1
\]

\textbf{2) Linearization Determinism Score $L_d(L)$}

JavaScript employs a prototype chain lookup algorithm, where property access is resolved deterministically by traversing the \texttt{[[Prototype]]} chain in a strictly defined order until a matching property is found or the chain terminates. This lookup procedure is formally specified in ECMAScript (e.g., via the \texttt{Get} and \texttt{[[Get]]} internal methods), ensuring a fully deterministic and unambiguous resolution order. Since JavaScript supports only single inheritance (one prototype link per object), there is no need for complex linearization strategies as in multiple inheritance systems, and the resolution order is uniquely defined by the chain structure.

\[
    L_d(\text{JavaScript}) = 1
\]

\textbf{3) Constraint Protection Coefficient $C_p(L)$}

$p_1=0$: the language does not provide \texttt{final} classes or methods. The inheritance via the \texttt{[[Prototype]]} chain remains dynamically extensible. 

$p_2=0$: ECMAScript lacks sealed or closed class hierarchies in the sense of restricting subclass sets (although \texttt{Object.seal} exists, it does not constrain inheritance relations). 

$p_3=0$: method overriding is implicit and no mandatory \texttt{override} annotation is specified. 

$p_4=0$: because all methods are effectively virtual by default due to dynamic dispatch through property lookup. 

$p_5=0$: since JavaScript does not enforce a formal separation between abstract classes and interfaces at the language level. 

$p_6=0$: because there are no compile-time rules prohibiting or restricting inheritance extensions, all checks are dynamic.

\[
    C_p(\text{JavaScript}) = \frac{0+0+0+0+0+0}{6} = 0,
\]

\textbf{4) Fragile Base Mitigation Factor $F_b(L)$}

$s_1=0$: the language does not define any mandatory explicit override annotation, method overriding is implicit and unchecked by the specification. 

$s_2=0$: JavaScript lacks strict method visibility controls at the specification level-while \texttt{\#private} fields exist, method visibility does not prevent overriding or exposure in inheritance hierarchies. 

$s_3=1$: ECMAScript explicitly requires invocation of \texttt{super()} in derived class constructors before accessing \texttt{this}, enforcing controlled superclass initialisation semantics. 

$s_4=0$: the language specification does not impose variance restrictions or static type constraints ensuring type-safe overriding. 

$s_5=0$: all methods are effectively overridable by default, with no built-in mechanism for non-overridable (final) methods. 

$s_6=0$: there are no compile-time override consistency checks in the dynamically typed specification. Therefore, the metric evaluates to

\[
    F_b(\text{JavaScript}) = \frac{1}{6} \approx 0.167,
\]

\textbf{Final metric value:}
\[
\begin{aligned} 
    ISRI(\text{JavaScript}) = 0.2 \cdot H_s(\text{JavaScript}) + 0.25 \cdot L_d(\text{JavaScript}) + 0.3 \cdot C_p(\text{JavaScript}) + 0.25 \cdot F_b(\text{JavaScript}) \\
    = 0.2 \cdot 1 + 0.25 \cdot 1 + 0.3 \cdot 0 + 0.25 \cdot \frac{1}{6} \\
    \approx 0.4917
\end{aligned}
\]

\subsubsection{Composition and Alternatives to Inheritance}

\textbf{1) Class-Based Delegation Score $D_{cls}(L)$}

JavaScript does not provide dedicated delegation keywords (such as \texttt{delegate}). However, delegation can be emulated through object composition, prototype chaining, or manual forwarding methods, which are idiomatic but not syntactically enforced constructs. 
The language does not include automatic method forwarding mechanisms-developers must explicitly define wrapper methods or use patterns (e.g., proxies) to delegate calls, which are not equivalent to built-in delegation features tied to class definitions

\[
    K_{del}(\text{JavaScript}) = 0.5, S_{fwd}(\text{JavaScript}) = 0 \Rightarrow D_{cls}(\text{JavaScript}) = 0.25
\]

\textbf{2) Trait Expressiveness Score $T_{exp}(L)$}

JavaScript resolves method name conflicts implicitly through property overwriting based on declaration or assignment order when mixin patterns (e.g., object spreading or prototype augmentation) are used. There are no formal mechanisms for explicit exclusion or aliasing, but neither does the language enforce compile-time errors. JavaScript mixin patterns cannot formally declare requirements (such as required methods or fields) within the language specification, any such constraints are informal and not enforced by the runtime or syntax.

\[
    R_{conf}(\text{JavaScript}) = 0, O_{mod}(\text{JavaScript}) = 0 \Rightarrow T_{exp}(\text{JavaScript}) = 0
\]

\textbf{3) Interface Composition Capacity $I_{comp}(L)$}

The language does not include native interface constructs or formally defined analogues at the specification level. 
There is no facility for declaring default method implementations, because of absence of interfaces. While JavaScript supports behavioural composition through mixins and object composition, these are not formal interface-based mechanisms and are not specified as such.

\[
    I_{comp}(\text{JavaScript}) = 0
\]

\textbf{Final metric value:}
\[
\begin{aligned}
    CRFI(\text{JavaScript}) = 0.3 \cdot D_{cls}(\text{JavaScript}) + 0.35 \cdot T_{exp}(\text{JavaScript}) + 0.35 \cdot I_{comp}(\text{JavaScript}) \\
    = 0.3 \cdot 0.25 + 0.35 \cdot 0 + 0.35 \cdot 0 = 0.075
\end{aligned}
\]

\subsubsection{Object Representation and Creation Model}

\textbf{1) Identity Semantics Score $S_{id}(L)$}

JavaScript allows the creation of objects via object literals and class-based syntax, both of which produce objects governed by reference semantics. All user-definable structured types are reference-based. Primitive types (e.g., Number, String, Boolean) are not user-definable in the specification sense and are excluded from $T_{user}(L)$.

\[
    S_{id}(\text{JavaScript}) = 1
\]

\textbf{2) Instantiation Paradigm Diversity $S_{cr}(L)$}

JavaScript provides: constructor-based instantiation via the \texttt{new} operator with functions or classes, object literal construction using \texttt{{}}, and prototype-based creation via \texttt{Object.create}, which directly specifies the prototype of a new object. Other patterns (e.g., factory functions) are reducible to these primitives and are not considered separate paradigms.

\[
    P_{supported}(\text{C\#}) = 4 \Rightarrow
    S_{cr}(\text{JavaScript}) = \frac{3}{4} = 0.75
\]

\textbf{3) Meta-level Extensibility Score $S_{meta}(L)$}

JavaScript provides: runtime structural modification of objects (e.g., adding, deleting, or redefining properties via \texttt{Object.defineProperty}) and, more importantly, the ability to override fundamental object interaction semantics through \texttt{Proxy} objects, which intercept operations such as property access, assignment, and even object construction. This constitutes modification not only of structure but also of the object interaction protocol itself, corresponding to creation and behaviour interception at a meta-level

\[
    S_{meta}(\text{JavaScript}) = \frac{3}{3}
\]

\textbf{Final metric value:}
\[
\begin{aligned}
    ORCI(\text{JavaScript}) = 0.35 \cdot S_{id}(\text{JavaScript}) + 0.3 \cdot S_{cr}(\text{JavaScript}) + 0.35 \cdot S_{meta}(\text{JavaScript}) \\
    = 0.35 \cdot 1 + 0.3 \cdot 0.75 + 0.35 \cdot 1 = 0.925
\end{aligned}
\]

\subsubsection{Polymorphism and Dispatch Mechanisms}

\textbf{1) Dispatch Scope Mechanism $M_{scope}(L)$}

JavaScript supports single dispatch through its object model, where method invocation is determined by the receiver object and resolved via the prototype chain.

$I_{single} = 1$

However, the language does not natively support multiple dispatch based on the runtime types of multiple arguments.

$I_{multi} = 0$ 

Similarly, predicate-based dispatch is not a built-in language feature-conditional logic can emulate it, but there is no formal dispatch mechanism based on predicates in the specification.

$I_{pred} = 0$

\[
    M_{scope}(\text{JavaScript}) = \frac{1}{5} = 0.2
\]

\textbf{2) Dispatch Safety Guarantee $S_{disp}(L)$}

The specification defines a dynamically typed language with no compile-time verification phase. The method dispatch correctness (e.g., property existence or callable status) is resolved at runtime via operations such as \texttt{[[Get]]} and \texttt{[[Call]]}, and failures result in runtime exceptions rather than static rejection. 

$I_{static}=0$

The language does not enforce exhaustive dispatch definitions, constructs such as \texttt{switch} statements or property-based dispatch do not require coverage of all possible cases, and missing branches are permitted without specification-level errors. Substituting into the metric,

$I_{exhaust}=0$

\[
    S_{disp}(\text{JavaScript}) = \frac{0 + 0}{2} = 0,
\]

\textbf{3) Virtualisation Optimization Potential $O_{virt}(L)$}

JavaScript does not support final or non-virtual methods, as all methods are dynamically dispatched and can be overridden.

$K_{final} = 0$

The language also lacks sealed class hierarchies that restrict subclassing.

$K_{sealed} = 0$

While JavaScript includes \texttt{static} methods in class syntax, these do not provide static dispatch guarantees or deVirtualisation hints in the specification sense, as method resolution remains dynamic.

$K_{static} = 0$

\[
    O_{virt}(\text{JavaScript}) = 0
\]

\textbf{Final metric value:}
\[
\begin{aligned}
    PDEI(\text{JavaScript}) = 0.35 \cdot M_{scope}(\text{JavaScript}) + 0.35 \cdot S_{disp}(\text{JavaScript}) + 0.3 \cdot O_{virt}(\text{JavaScript}) \\
    = 0.35 \cdot 0.2 + 0.35 \cdot 0 + 0.3 \cdot 0 = 0.07
\end{aligned}
\]

\subsubsection{State Model, Mutability, and Concurrency Guarantees}

\textbf{1) Default Mutability Score $M_{def}(L)$}

All standard objects (including arrays and functions) are mutable unless explicitly restricted via mechanisms such as \texttt{Object.freeze} or by using immutable patterns at the application level. Therefore, mutability does not require explicit annotation, immutability must be enforced explicitly by the programmer.

\[
    M_{def}(\text{JavaScript}) = 0
\]

\textbf{2) Static Reference Constraint Density $S_{dens}(L)$}

First, constructs that introduce or manipulate references: variable bindings (\texttt{var}, \texttt{let}, \texttt{const}), object property access, array indexing, and function parameter passing, all of which operate on references to objects in the heap via the specification’s internal Reference Record model. 

$N_{ref}(\text{JavaScript}) = 4$

Second, $N_{constr}(\text{JavaScript})$ counts only those constructs that enforce static (compile-time) constraints on reference usage. JavaScript, being a dynamically typed language, does not impose compile-time aliasing or mutability constraints.

\[
    S_{dens}(\text{JavaScript}) = 0
\]

\textbf{3) Concurrency Primitive Safety $P_{conc}(L)$}

First, JavaScript uses a single-threaded event loop, which reduces the risk of data races. However, with the advent of \texttt{SharedArrayBuffer} and \texttt{Atomics}, parallel work with shared memory has become possible, where correctness depends on the programmer, and not on the built-in guarantees of the language.

$S_{race}(\text{JavaScript}) = 0$

Second, JavaScript provides high-level concurrency primitives such as \texttt{Promise}, \texttt{async/await}, and event-driven callbacks, which correspond to structured concurrency (task-based asynchronous computation) rather than actor-style isolation or CSP channels.

$S_{model}(\text{JavaScript}) = 0.5$. 

\[
    P_{conc}(\text{JavaScript}) = \frac{0 + 0.5}{2} = 0.25
\]

\textbf{Final metric value:}
\[
\begin{aligned}
    SCSI(\text{JavaScript}) = 0.35 \cdot M_{def}(\text{JavaScript}) + 0.25 \cdot S_{dens}(\text{JavaScript}) + 0.4 \cdot P_{conc}(\text{JavaScript}) \\
    = 0.35 \cdot 0 + 0.25 \cdot 0 + 0.4 \cdot 0.25 = 0.1
\end{aligned}
\]

\subsubsection{Support for Formal Specification and Verifiability}

\textbf{1) Native Contract Constructs Score $C_{nat}(L)$}

JavaScript does not provide first-class language constructs for specifying preconditions, postconditions, or invariants, such checks are implemented manually using conditional statements or via external libraries (e.g., assertion frameworks), which are not part of the core language semantics. Specifically, there are no native syntactic forms analogous to \texttt{require}, \texttt{ensure}, or invariant clauses within class or function definitions.

\[
    C_{nat}(\text{JavaScript}) = \frac{0 + 0 + 0}{3} = 0
\]

\textbf{2) Verification Enforcement Level $V_{enf}(L)$}

JavaScript does not support static validation at compile time due to dynamic typing. Errors are detected only at runtime via type checking, exceptions, and built-in errors (e.g. \texttt{TypeError}).

\[
    V_{enf}(\text{JavaScript}) = 0.5
\]

\textbf{Final metric value:}
\[
\begin{aligned}
    FSVI(\text{JavaScript}) = 0.45 \cdot C_{nat}(\text{JavaScript}) + 0.55 \cdot V_{enf}(\text{JavaScript}) \\
    = 0.45 \cdot 0 + 0.55 \cdot 0.5 = 0.275
\end{aligned}
\]


\subsection{Metrics Evaluation for Python}

Python is a dynamically typed object-oriented language with flexible architecture and adaptive code execution. The philosophy of the language emphasises simplicity and flexibility, instead of strict guarantees of type safety and formal verification. That is why Python provides extensive OOP capabilities, but less support for strict typing, encapsulation, and verification. In this section, metrics for Python are calculated based on the official specification of the language and its documented semantics.

\subsubsection{1) Type System and Subtyping Model}

\textbf{1) Nominal–Structural Typing Score $S_{nom}(L)$}

In the case of Python, nominal typing is clearly supported through its class system, where subtype relationships are explicitly declared via inheritance \cite{pep544}. 

$I_{nom}(\text{Python})=1$

Structural typing is also supported through the introduction of \textit{Protocols} \cite{pep544}, which define subtyping based on the presence of methods and properties rather than explicit inheritance.

$I_{struct}(\text{Python})=1$ .

\[
    S_{nom}(\text{Python}) = 1
\]

\textbf{2) Variance Formalization Score $S_{var}(L)$}

In Python, variance is supported in the type hinting system \cite{pep484} through the use of generic type variables (e.g., \texttt{TypeVar}) with explicit annotations such as \texttt{covariant=True} and \texttt{contravariant=True}, hence co-, contra-, and invariance can be directly declared.

Python does not provide strict type safety at runtime and uses external validation tools (e.g. mypy). There is no formally proven model of typical security in the specification, so the language applies only practical constraints to reduce errors.

\[
\begin{aligned}
    V_{exp}(\text{Python}) = 1, V_{sound}(\text{Python}) = 0.5 \\
    S_{var}(\text{Python}) = 0.75
\end{aligned}
\]

\textbf{3) Inheritance–Subtyping Separation Score $S_{rel}(L)$}

First, Python supports explicit subtyping without inheritance via structural typing mechanisms such as \texttt{typing.Protocol}, allowing a class to be considered a subtype based on method signatures rather than inheritance hierarchy.

$I_{dec}(\text{Python})=1$

Second, subtyping rules in Python’s static type system are defined independently of inheritance semantics: nominal inheritance (class-based) and structural subtyping (protocol-based) coexist, and type compatibility is determined by type checkers without requiring subclass relationships.

$I_{ind}(\text{Python})=1$

Third, although the distinction between inheritance and subtyping is clearly articulated in the typing PEPs, it is not fully formalized in a rigorous mathematical specification within the core language definition, remaining partly informal and tool-dependent. Thus, substituting these values into the formula yields.

$I_{form}(\text{Python})=0.5$

\[
    S_{rel}(\text{Python}) = \frac{1 + 1 + 0.5}{3} = \frac{2.5}{3} \approx 0.83,
\]

\textbf{Final metric value:}
\[
\begin{aligned}
    TSSI(\text{Python}) = 0.3 \cdot S_{nom}(\text{Python}) + 0.3 \cdot S_{var}(\text{Python}) + 0.4 \cdot S_{rel}(\text{Python}) \\
    = 0.3 \cdot 1 + 0.3 \cdot 0.75 + 0.4 \cdot \frac{5}{6} \\
    \approx 0.8583
\end{aligned}
\]

\subsubsection{Abstraction and Encapsulation Mechanisms}

\textbf{1) Visibility Granularity Index $V(L)$}

In Python, access control is not enforced through strict modifiers but follows a convention-based model as described in the language reference: 

\begin{itemize}
    \item public members (default)
    \item protected members indicated by a single leading underscore (e.g., \texttt{\_x})
    \item private members using name mangling with double leading underscores (e.g., \texttt{\_\_x})
\end{itemize}

Although these mechanisms do not enforce strict encapsulation at the compiler or runtime level, they are semantically distinct and consistently recognized by the specification and tooling. 


\[
\begin{aligned}
    n_{levels}(\text{Python}) = 3, n_{levels}(\text{C\#})=7 \\
    V(\text{Python})=0.\frac{3}{7}
\end{aligned}
\]


\textbf{2) Module Isolation Strength $M(L)$}

In Python, the primary modularisation construct is the module itself (i.e., a file), along with packages that group modules hierarchically.

$b_{total}(\text{Python}) = 2$ (modules and packages)

Python does not provide strict export controls: the \texttt{\_\_all\_\_} variable affects only imports via \texttt{from module import *}, but does not prohibit direct access to other attributes. All names remain accessible, reflecting the principle of open access instead of strict encapsulation.

$b_{strict}(\text{Python}) = 0$

\[
    M(\text{Python})=0.5
\]


\textbf{3) Interface Abstraction Capacity $I(L)$}

In Python, abstraction is supported through abstract base classes from the \texttt{abc} module and protocols (\textit{Protocols}) that describe behavior independent of inheritance.

$n_{abstraction\_constructs}(\text{Python}) = 2$

The primary object-defining construct in Python is the class. Functions are not object-defining in the same structural sense.

$n_{object}(\text{Python}) = 1$.

\[
    I(\text{Python})=\frac{2}{1}=2 \Rightarrow \text{normalized} \approx 1
\]


\textbf{4) Encapsulation Enforcement Coefficient $E(L)$}

In Python, encapsulation is not strictly enforced: all object attributes are accessible via standard mechanisms such as \texttt{getattr}, \texttt{setattr}, direct dictionary manipulation through \texttt{\_\_dict\_\_}, and introspection/reflection facilities defined in the language specification. Python also lacks restrictions on reflective capabilities and does not provide a type system that enforces access boundaries at runtime. Therefore, encapsulation violations are not prevented, and the language explicitly permits unrestricted access to object internals.

\[
    E(\text{Python}) = 0.2
\]

\textbf{Final metric value:}
\[
\begin{aligned}
    AE(\text{Python}) = 0.2 \cdot V(\text{Python}) + 0.25 \cdot M(\text{Python}) + 0.25 \cdot I(\text{Python}) + 0.3 \cdot E(\text{Python}) \\
    = 0.2 \cdot \frac{3}{7} + 0.25 \cdot 0.5 + 0.25 \cdot 1 + 0.3 \cdot 0.2 \\
    \approx 0.5207
\end{aligned}
\]

\subsubsection{Inheritance Semantics and Hierarchy Management}

\textbf{1) Hierarchy Simplicity Factor $H_s(L)$}

In Python, the class definition syntax explicitly allows multiple inheritance, i.e., a class can inherit from an arbitrary number of base classes (e.g., \texttt{class C(A, B, ...)}), implying that $M(\text{Python}) > 1$. Since there is no fixed upper bound in the specification, Python effectively reaches the maximal considered value in comparative analysis.
\[
    M(\text{Python}) = M_{max}
    \Rightarrow H_s(\text{Python}) = 1 - \frac{M(\text{Python}) - 1}{M_{max} - 1} = 0
\]


\textbf{2) Linearization Determinism Score $L_d(L)$}

In Python, the MRO is explicitly defined in the language specification via the C3 linearization algorithm, which ensures a consistent and monotonic resolution order for multiple inheritance hierarchies. This algorithm is formally described and guarantees determinism across all compliant implementations, eliminating ambiguity in method lookup. Since Python fully specifies its linearization strategy and enforces it uniformly, it satisfies the condition for a formally defined and deterministic MRO.

\[
    L_d(\text{Python})=1
\]


\textbf{3) Constraint Protection Coefficient $C_p(L)$}

$p_1=0$: \texttt{typing.final} is not enforced by the core language semantics and remains advisory for static type checkers only.

$p_2=0$: Python does not provide sealed or closed class hierarchies.

$p_3=0$: method overriding does not require an explicit \texttt{override} annotation.

$p_4=0$: all methods are virtual by default due to dynamic dispatch.

$p_5=1$: Python distinguishes abstract base classes through the \texttt{abc} module, checking the implementation of abstractions when instantiating objects.

$p_6=0$: there are no compile-time restrictions on inheritance in the dynamically typed core language.

\[
    C_p(\text{Python}) = \frac{0+0+0+0+1+0}{6} \approx 0.17,
\]


\textbf{4) Fragile Base Mitigation Factor $F_b(L)$}

$s_1=0$: Python does not require the required \texttt{override}.

$s_2=0$: access control is based on conventions, not strict modifiers.

$s_3=0$: calling \texttt{super()} is optional, and the method resolution order allows for implicit behaviour.

$s_4=1$: the typing system supports the rules of variation, checked by type checkers.

$s_5=0$: all methods are virtual by default and can be overridden.

$s_6=0$: there are no built-in compile-time guarantees; checks are performed by external tools.

\[
    F_b(\text{Python}) = \frac{0+0+0+1+0+0}{6} \approx 0.17,
\]

\textbf{Final metric value:}
\[
\begin{aligned}
    ISRI(\text{Python}) = 0.2 \cdot H_s(\text{Python}) + 0.25 \cdot L_d(\text{Python}) + 0.3 \cdot C_p(\text{Python}) + 0.25 \cdot F_b(\text{Python}) \\
    = 0.2 \cdot 0 + 0.25 \cdot 1 + 0.3 \cdot \frac{1}{6} + 0.25 \cdot \frac{1}{6} \\
    \approx 0.3417
\end{aligned}
\]

\subsubsection{Composition and Alternatives to Inheritance}

\textbf{1) Class-Based Delegation Score $D_{cls}(L)$}

Python does not provide built-in mechanisms or keywords for delegation. Developer should implement delegation manually using composition, intermediary methods, or dynamic methods.

\[
    K_{del}(\text{Python})=0.5, S_{fwd}(\text{Python})=0 \Rightarrow D_{cls}(\text{Python})=0.25
\]

\textbf{2) Trait Expressiveness Score $T_{exp}(L)$}

Python implement mixins through multiple inheritance, and resolve naming conflicts deterministically using MRO (C3 linearization). The language does not provide special tools for declaring required methods or properties in mixins.

\[
    R_{conf}(\text{Python}) = 0.5, O_{conf}(\text{Python}) = 0 \Rightarrow T_{exp}=0.25
\]

\textbf{3) Interface Composition Capacity $I_{comp}(L)$}

In Python, there is no dedicated \textit{interface} construct in the core language specification. However, abstract base classes (ABCs) and \textit{Protocols} (PEP 544) serve as functional analogues. A class may inherit from multiple abstract base classes or conform to multiple protocols without restriction. 
Furthermore, abstract base classes in Python may include concrete method implementations alongside abstract methods, enabling default behaviour reuse, while protocols may also define method bodies.

\[
    N_{impl}(\text{Python}) = 1, D_{def}(\text{Python}) = 1 \Rightarrow I_{comp} = 1
\]

\textbf{Final metric value:}
\[
\begin{aligned}
    CRFI(\text{Python}) = 0.3 \cdot D_{cls}(\text{Python}) + 0.35 \cdot T_{exp}(\text{Python}) + 0.35 \cdot I_{comp}(\text{Python}) \\
    = 0.3 \cdot 0.25 + 0.35 \cdot 0.25 + 0.35 \cdot 1 \\
    = 0.5125
\end{aligned}
\]

\subsubsection{Object Representation and Creation Model}

\textbf{1) Identity Semantics Score $S_{id}(L)$}

In Python, the language specification defines a uniform object model in which all user-defined types (i.e., classes) produce objects with identity, accessible via operations such as \texttt{id()} and compared using the \texttt{is} operator. There are no alternative user-definable value-type categories with copy semantics distinct from reference behaviour, even immutable types preserve identity at the object level.
All user-defined types are reference types:

\[
    S_{id}(\text{Python})=1
\]

\textbf{2) Instantiation Paradigm Diversity $S_{cr}(L)$}

The object creation is primarily class-based via constructor invocation (e.g., calling a class object, which internally triggers \texttt{\_\_new\_\_} and \texttt{\_\_init\_\_}), and literal-based construction for built-in container types (e.g., lists, dictionaries, sets), both formally defined in the language. Additionally, Python supports factory-based instantiation patterns through first-class functions and metaclass customization, but these are extensions of the class-based model rather than distinct paradigms such as prototype cloning or actor spawning.
\[
    P_{supported}(\text{C\#}) = 4 \Rightarrow
    S_{cr}(\text{Python}) = \frac{2}{4}
\]

\textbf{3) Meta-level Extensibility Score $S_{meta}(L)$}

Python provides extensive meta-programming capabilities: it supports full introspection (e.g., via \texttt{dir()}, \texttt{type()}, reflection APIs), runtime structural modification (e.g., dynamic attribute addition, modification of \texttt{\_\_dict\_\_}), and, critically, allows overriding the object creation protocol through metaobjects such as metaclasses and customization hooks like \texttt{\_\_new\_\_}, \texttt{\_\_init\_\_}, and \texttt{\_\_call\_\_}.

\[
    S_{meta}(\text{Python})= \frac{3}{3}
\]

\textbf{Final metric value:}
\[
\begin{aligned}
    ORCI(\text{Python}) = 0.35 \cdot S_{id}(\text{Python}) + 0.3 \cdot S_{cr}(\text{Python}) + 0.35 \cdot S_{meta}(\text{Python}) \\
    = 0.35 \cdot 1 + 0.3 \cdot \frac{2}{4} + 0.35 \cdot 1 \\
    = 0.85
\end{aligned}
\]

\subsubsection{Polymorphism and Dispatch Mechanisms}

\textbf{1) Dispatch Scope Mechanism $M_{scope}(L)$}

In Python, single dispatch is natively supported through method invocation on the receiver object, as defined by the object model and method resolution order.

$I_{single}(\text{Python}) = 1$

The language does not support multiple dispatch directly. Libraries can overload functions with only one argument, but full-fledged dispatching for all parameters is not included in the specification.

$I_{multi}(\text{Python}) = 0$

The language does not directly support predicate-based dispatch. The conditions need to be written in the code manually, and not through the dispatch rules.

$I_{pred}(\text{Python}) = 0$

\[
    M_{scope}(\text{Python}) = \frac{1}{5} = 0.2
\]

\textbf{2) Dispatch Safety Guarantee $S_{disp}(L)$}

First, Python does not provide compile-time guarantees of dispatch correctness: method resolution and function dispatch are determined at runtime via dynamic typing and late binding, while static type checkers (e.g., \texttt{mypy}) remain external and non-mandatory. 

$I_{static}(\text{Python})=0$

Second, the language does not enforce exhaustiveness of dispatch definitions, even with structural pattern matching (\texttt{match} statement, PEP 634), exhaustiveness checking is not required or guaranteed by the specification. 

$I_{exhaust}(\text{Python})=0$

\[
    S_{disp}(\text{Python}) = \frac{0 + 0}{2} = 0,
\]

\textbf{3) Virtualisation Optimization Potential $O_{virt}(L)$}

In Python, the specification does not provide enforced \texttt{final} or non-virtual methods. All methods are dynamically dispatched and overrideable, and annotations such as \texttt{final} are not enforced at runtime.

$K_{final}(\text{Python}) = 0$

Similarly, Python does not support sealed class hierarchies that restrict subclassing.

$K_{sealed}(\text{Python}) = 0$ 

Although Python includes \texttt{@staticmethod}, this construct does not serve as a static dispatch optimization hint but rather defines methods without implicit receiver binding, and does not eliminate dynamic lookup semantics.

$K_{static}(\text{Python}) = 0$

\[
    O_{virt}(\text{Python})=0
\]

\textbf{Final metric value:}
\[
\begin{aligned}
    PDEI(\text{Python}) = 0.35 \cdot M_{scope}(\text{Python}) + 0.35 \cdot S_{disp}(\text{Python}) + 0.3 \cdot O_{virt}(\text{Python}) \\
    = 0.35 \cdot 0.2 + 0.35 \cdot 0 + 0.3 \cdot 0 \\
    = 0.07
\end{aligned}
\]

\subsubsection{State Model, Mutability, and Concurrency Guarantees}

\textbf{1) Default Mutability Score $M_{def}(L)$}

In Python, according to its language specification and data model, objects are not uniformly immutable by default: core built-in types such as integers, floats, strings, and tuples are immutable, whereas lists, dictionaries, and sets are mutable without requiring any explicit annotation. The language does not enforce immutability as a default semantic constraint, nor does it require developers to explicitly declare objects as mutable, mutability is an inherent property of each type and is directly accessible without syntactic modifiers.

\[
    M_{def}(\text{Python})=0
\]

\textbf{2) Static Reference Constraint Density $S_{dens}(L)$}

The type system is checked at runtime, not at compilation. The language does not limit references, mutability, and data ownership at the type level. There are type hints, but they are not strict compile-time constraints.

\[
    S_{dens}(\text{Python}) = 0
\]

\textbf{3) Concurrency Primitive Safety $P_{conc}(L)$}

First, Python does not provide compile-time guarantees of data race freedom in its type system. It relies primarily on programmer discipline when using shared-memory concurrency (e.g., via the \texttt{threading} module). Although the Global Interpreter Lock (GIL) in CPython restricts true parallel execution of bytecode, it does not eliminate logical data races on shared mutable state, nor is it a formal runtime race detection mechanism as defined by the metric. 

$S_{race}(\text{Python}) = 0$ 

Second, Python provides structured concurrency abstractions such as \texttt{asyncio} with tasks, futures, and event loops, which correspond to the “structured concurrency” category rather than purely low-level primitives. However, it does not natively enforce higher-level models like Actors or CSP as first-class language constructs. 

$S_{model}(\text{Python}) = 0.5$

\[
    P_{conc}(\text{Python})=0.25
\]

\textbf{Final metric value:}
\[
\begin{aligned}
    SCSI(\text{Python}) = 0.35 \cdot M_{def}(\text{Python}) + 0.25 \cdot S_{dens}(\text{Python}) + 0.4 \cdot P_{conc}(\text{Python}) \\
    = 0.35 \cdot 0 + 0.25 \cdot 0 + 0.4 \cdot 0.25 = 0.1
\end{aligned}
\]

\subsubsection{Support for Formal Specification and Verifiability}

\textbf{1) Native Contract Constructs Score $C_{nat}(L)$}

Python does not have built-in precondition support. Developer can implement preconditions using \texttt{assert} or third-party libraries

$n_{pre}(\text{Python}) = 0$ 

Python lacks native postcondition constructs. Only assertion and external tools may be used, but ther are not part of the language.

$n_{post}(\text{Python}) = 0$ 

Python does not provide first-class invariant constructs at the language level, they must be manually encoded or delegated to third-party libraries.

$n_{inv}(\text{Python}) = 0$

\[
    C_{nat}(\text{Python}) = 0
\]

\textbf{2) Verification Enforcement Level $V_{enf}(L)$}

Python does not support static validation and proof of correctness during compilation. Type hints are optional and work only through external analyzers. Checks are possible at runtime, for example through `assert` and dynamic type checking.

\[
    V_{enf}(\text{Python})=0.5
\]

\textbf{Final metric value:}
\[
\begin{aligned}
    FSVI(\text{Python}) = 0.45 \cdot C_{nat}(\text{Python}) + 0.55 \cdot V_{enf}(\text{Python}) \\
    = 0.45 \cdot 0 + 0.55 \cdot 0.5 = 0.275
\end{aligned}
\]

\subsection{Metrics Evaluation for Ruby}

Ruby is a dynamically typed, purely object-oriented programming language with a highly reflective metaobject protocol and a strong emphasis on programmer flexibility. According to Flanagan and Matsumoto, “Ruby is a pure object-oriented language: all data in Ruby are objects” \cite{flanagan2008ruby}. However, this flexibility comes at the cost of formal guarantees in several aspects of object modeling, particularly in type safety, verification, and concurrency. In this subsection, we systematically evaluate Ruby using the previously defined metrics, grounding each numerical assignment in the formal language specification and authoritative literature such as \cite{flanagan2008ruby}.

\subsubsection{1) Type System and Subtyping Model}

Ruby employs a dynamically typed system without a formal static type specification. Subtyping is implicit and primarily behavioural (duck typing), rather than formally defined.

\textbf{1) Nominal–Structural Typing Score $S_{nom}(L)$}

Ruby is a dynamically typed, object-oriented language that employs \emph{duck typing}, meaning that type compatibility is determined by the presence of methods rather than explicit type declarations. However, this behaviour is not formalized as structural subtyping with explicit typing judgments or inference rules in the specification.

$I_{struct}(\text{Ruby})=0$

At the same time, Ruby supports nominal typing through its class-based object model, where subtyping relationships are explicitly declared via inheritance (e.g., using \texttt{<} for subclassing), and method lookup follows the class hierarchy, which satisfies the condition for nominal typing.

$I_{nom}(\text{Ruby})=1$

\[
    S_{nom}(\text{Ruby}) = \frac{1 + 0}{2} = 0.5
\]

\textbf{2) Variance Formalization Score $S_{var}(L)$}

Ruby does not natively support parametric polymorphism (generics) in its core language specification. Although recent extensions such as RBS and type checkers (e.g., Sorbet) introduce generic-like constructs, these are external and not part of the official language semantics.

$V_{exp}(\text{Ruby})=0$

The variance (covariance/contravariance) is not formally defined within the Ruby language specification and there are no typing judgments or proofs ensuring type safety for such mechanisms, the soundness component is also absent.

$V_{sound}(\text{Ruby})=0$

\[
    S_{var}(\text{Ruby}) = 0
\]

\textbf{3) Inheritance–Subtyping Separation Score $S_{rel}(L)$}

First, Ruby does not provide explicit constructs for subtyping independent of inheritance (e.g., no interfaces or type declarations distinct from class hierarchies), and type compatibility is determined dynamically via method availability rather than formally declared subtype relations.

$I_{dec}(\text{Ruby})=0$

Second, subtyping is not defined as an independent formal relation in the specification but is effectively derived from class inheritance and mixin inclusion semantics.

$I_{ind}(\text{Ruby})=0$

Third, the Ruby specification does not formally distinguish inheritance and subtyping as separate relations with independent typing judgments.

$I_{form}(\text{Ruby})=0$.

\[
    S_{rel}(\text{Ruby})=\frac{0+0+0}{3}=0
\]

\textbf{Final metric value:}

\[
\begin{aligned}
    TSSI(\text{Ruby}) = 0.3 \cdot S_{nom}(\text{Ruby}) + 0.3 \cdot S_{var}(\text{Ruby}) + 0.4 \cdot S_{rel}(\text{Ruby}) \\
    = 0.3 \cdot 0.5 + 0.3 \cdot 0 + 0.4 \cdot 0 = 0.15
\end{aligned}
\]

\subsubsection{Abstraction and Encapsulation Mechanisms}

\textbf{1) Visibility Granularity Index $V(L)$}

Ruby provides three semantically distinct visibility levels for methods: \texttt{public}, \texttt{protected}, and \texttt{private} \cite{flanagan2008ruby}.

\[
    n_{levels}(\text{C\#})=7 \Rightarrow V(\text{Ruby}) = \frac{3}{7}
\]

\textbf{2) Module Isolation Strength $M(L)$}

Ruby provides modularisation through \texttt{class} and \texttt{module} constructs, which define namespaces and method lookup scopes. 

$b_{total}(\text{Ruby})=2$

However, Ruby does not support explicit export control lists or enforced visibility at the module or file boundary level, and all constants within a module or class are accessible via namespace qualification.

$b_{strict}(\text{Ruby})=0$

\[
    M(\text{Ruby}) = 0
\]

\textbf{3) Interface Abstraction Capacity $I(L)$}

Ruby does not provide dedicated interface or protocol constructs. It relies on implicit behavioural abstraction via duck typing and mixins (modules), where \texttt{module} can define method signatures but also includes implementations, thus not serving as a purely abstract specification mechanism. Therefore, only one abstraction-related construct (\texttt{module}) can be partially attributed to behavioural abstraction. 

$n_{abstraction\_constructs}(\text{Ruby})=1$

The language defines two primary object-defining constructs, \texttt{class} and \texttt{module}, both of which can produce objects or object-like namespaces.

$n_{object}(\text{Ruby})=2$

\[
    I(\text{Ruby}) = \frac{1}{2} = 0.5
\]

\textbf{4) Encapsulation Enforcement Coefficient $E(L)$}

Ruby allows users to call private methods and change the internal state of objects through reflection and metaprogramming mechanisms, such as \texttt{send}, \texttt{instance\_eval} and \texttt{class\_eval}. Developers can dynamically read and modify object variables without checking at the compilation stage. Ruby also does not strictly control the visibility of methods at runtime when using reflection. Therefore, the language allows for a violation of encapsulation.

\[
    E(\text{Ruby}) = 0
\]

\textbf{Final metric value:}

\[
\begin{aligned}
    AE(\text{Ruby}) = 0.2 \cdot V(\text{Ruby}) + 0.25 \cdot M(\text{Ruby}) + 0.25 \cdot I(\text{Ruby}) + 0.3 \cdot E(\text{Ruby}) \\
    = 0.2 \cdot \frac{3}{7} + 0.25 \cdot 0 + 0.25 \cdot 0.5 + 0.3 \cdot 0 \\
    \approx 0.2107
\end{aligned}
\]

\subsubsection{Inheritance Semantics and Hierarchy Management}

\textbf{1) Hierarchy Simplicity Factor $H_s(L)$}

Ruby enforces single inheritance for classes, meaning each class can have exactly one direct superclass (specified using the \texttt{<} syntax), while additional code reuse is achieved via mixins (\texttt{module} inclusion) rather than multiple inheritance. The formal inheritance hierarchy remains strictly single-parent.

\[
    H_s(\text{Ruby}) = 1
\]

\textbf{2) Linearization Determinism Score $L_d(L)$}

Ruby specifies a well-defined and deterministic method lookup algorithm based on a linearized ancestor chain, where the search order follows the class, included modules (in reverse order of inclusion), superclass, and so on up to \texttt{BasicObject}. This linearization is explicitly defined in the language semantics and consistently applied across implementations, ensuring predictable method dispatch even in the presence of mixins \cite{flanagan2008ruby}.

\[
    L_d(\text{Ruby}) = 1
\]

\textbf{3) Constraint Protection Coefficient $C_p(L)$}

$p_1=0$: Ruby does not support final classes and final methods.

$p_2=0$: Ruby does not support closed hierarchies.

$p_3=0$: Ruby does not require explicit use of \texttt{override} annotations.

$p_4=0$: Ruby initiate methods virtual by default.

$p_5=0$: Ruby does not separate abstractions and interfaces because the language does not support formal interfaces.

$p_6=0$: Ruby allows developers to freely extend and redefine classes at runtime.

\[
    C_p(\text{Ruby})=\frac{0+0+0+0+0+0}{6}=0
\]

\textbf{4) Fragile Base Mitigation Factor $F_b(L)$}

$s_1=0$: Ruby does not require explicit use of \texttt{override} annotations. 

$s_2=1$: Ruby supports the visibility modifiers (\texttt{public}, \texttt{protected}, \texttt{private}).

$s_3=0$: Ruby does not require superclass invocation. 

$s_4=0$: Ruby does not support variation constraints due to the lack of static typing.

$s_5=0$: Ruby allows developers to override default methods.

$s_6=0$: Ruby does not check for redefinition during compilation.

\[
    F_b(\text{Ruby})=\frac{0+1+0+0+0+0}{6}=\frac{1}{6}\approx 0.17
\]

\textbf{Final metric value:}

\[
\begin{aligned}
    ISRI(\text{Ruby}) = 0.2 \cdot H_s(\text{Ruby}) + 0.25 \cdot L_d(\text{Ruby}) + 0.3 \cdot C_p(\text{Ruby}) + 0.25 \cdot F_b(\text{Ruby}) \\
    = 0.2 \cdot 1 + 0.25 \cdot 1 + 0.3 \cdot 0 + 0.25 \cdot \frac{1}{6} \\
    \approx 0.4917
\end{aligned}
\]

\subsubsection{Composition and Alternatives to Inheritance}

\textbf{1) Class-Based Delegation Score $D_{cls}(L)$}

Ruby does not support built-in delegation keywords. The language allows you to implement delegation through \texttt{Forwardable}, \texttt{SimpleDelegator}, \texttt{method\_missing} and metaprogramming mechanisms. These tools automatically redirect method calls to the internal object and reduce the amount of boilerplate code.

\[
    K_{del}(\text{Ruby})=0.5, S_{fwd}(\text{Ruby})=0 \Rightarrow D_{cls}(\text{Ruby})=0.25
\]

\textbf{2) Trait Expressiveness Score $T_{exp}(L)$}

Ruby resolves method conflicts in mixin modules through the connection order: the last connected module gets priority. The language does not require explicit exclusion or renaming of methods. Ruby also does not allow modules to set mandatory methods or constraints for classes, so mixin modules rely on implicit assumptions and runtime checks.

\[
    R_{conf}(\text{Ruby})=0.5, O_{mod}(\text{Ruby})=0 \Rightarrow T_{exp}(\text{Ruby})=0.25
\]

\textbf{3) Interface Composition Capacity $I_{comp}(L)$}

Ruby does not define interfaces as separate constructs. However, modules serve as compositional units that can be included in classes without restriction on their number, effectively acting as interface-like mechanisms. 
Furthermore, unlike traditional interfaces, Ruby modules can contain full method implementations, allowing them to act as reusable behavioural units rather than pure type specifications

\[
    N_{iml}(\text{Ruby})=1, D_{def}(\text{Ruby})=1 \Rightarrow I_{comp}(\text{Ruby}) = 1
\]

\textbf{Final metric value:}

\[
\begin{aligned}
    CRFI(\text{Ruby}) = 0.3 \cdot D_{cls}(\text{Ruby}) + 0.35 \cdot T_{exp}(\text{Ruby}) + 0.35 \cdot I_{comp}(\text{Ruby}) \\
    = 0.3 \cdot 0.25 + 0.35 \cdot 0.25 + 0.35 \cdot 1 \\
    = 0.5125
\end{aligned}
\]

\subsubsection{Object Representation and Creation Model}

\textbf{1) Identity Semantics Score $S_{id}(L)$}

Ruby follows a uniform object model in which all user-defined entities are objects created via \texttt{class} or \texttt{module}, and all variables store references to these objects rather than values

\[
    S_{id}(\text{Ruby}) = 1
\]

\textbf{2) Instantiation Paradigm Diversity $S_{cr}(L)$}

Ruby primarily supports class-based instantiation via the \texttt{new} method (which internally invokes \texttt{allocate} and \texttt{initialize}), and also provides literal-based construction for built-in object types such as arrays, hashes, strings, and numeric values, which constitute a second paradigm of object creation. The language supports advanced metaprogramming facilities such as \texttt{Class.new}, \texttt{define\_method}, and dynamic evaluation via \texttt{eval}, which enable runtime generation of classes and objects. Ruby does not natively support prototype-based cloning as a primary paradigm (despite \texttt{clone} and \texttt{dup}, which operate on existing objects rather than serving as independent instantiation models), nor actor-based spawning in its core specification.

\[
    P_{supported}(\text{C\#}) = 4 \Rightarrow S_{cr}(\text{Ruby}) = \frac{3}{4}
\]

\textbf{3) Meta-level Extensibility Score $S_{meta}(L)$}

Ruby provides full reflective capabilities, including introspection (e.g., \texttt{methods}, \texttt{instance\_variables}), structural modification (e.g., \texttt{define\_method}, \texttt{class\_eval}, \texttt{module\_eval}), and, critically, the ability to override the object creation protocol via hooks such as \texttt{allocate}, \texttt{initialize}, and dynamic class construction with \texttt{Class.new}. These features allow programmers to intercept and redefine how objects are instantiated and structured at runtime, satisfying the highest level of intercession.

\[
    S_{meta}(\text{Ruby}) = \frac{3}{3}
\]

\textbf{Final metric value:}

\[
\begin{aligned}
    ORCI(\text{Ruby}) = 0.35 \cdot S_{id}(\text{Ruby}) + 0.3 \cdot S_{cr}(\text{Ruby}) + 0.35 \cdot S_{meta}(\text{Ruby}) \\
    = 0.35 \cdot 1 + 0.3 \cdot \frac{3}{4} + 0.35 \cdot 1 \\
    = 0.925
\end{aligned}
\]

\subsubsection{Polymorphism and Dispatch Mechanisms}

\textbf{1) Dispatch Scope Mechanism $M_{scope}(L)$}

Ruby natively supports single dispatch, where method selection is based on the receiver object and resolved through a deterministic method lookup chain.

$I_{single}(\text{Ruby})=1$

Ruby does not provide native multiple dispatch based on the dynamic types of multiple arguments.

$I_{multi}(\text{Ruby})=0$

The predicate-based dispatch is not part of the formal method dispatch mechanism in the specification. Although conditional logic can be implemented manually within methods, it is not integrated into the dispatch system itself.

$I_{pred}(\text{Ruby})=0$

\[
    M_{scope}(\text{Ruby}) = \frac{1}{5} = 0.2
\]

\textbf{2) Dispatch Safety Guarantee $S_{disp}(L)$}

First, Ruby is dynamically typed and does not perform compile-time checking of method existence or dispatch correctness, method resolution occurs at runtime and may raise exceptions such as \texttt{NoMethodError}.

$I_{static}(\text{Ruby})=0$ 

Second, Ruby does not enforce exhaustive dispatch definitions (e.g., no requirement to cover all possible cases in polymorphic behaviour or pattern matching at the specification level).

$I_{exhaust}(\text{Ruby})=0$

\[
    S_{disp}(\text{Ruby})=\frac{0+0}{2}=0
\]

\textbf{3) Virtualisation Optimization Potential $O_{virt}(L)$}

Ruby does not support \texttt{final} or non-virtual methods, as all methods are dynamically dispatched and can be overridden at runtime.

$K_{final}(\text{Ruby})=0$

Similarly, Ruby does not provide sealed class hierarchies or any mechanism to restrict subclassing.

$K_{sealed}(\text{Ruby})=0$

Furthermore, the language lacks static dispatch constructs or annotations (such as \texttt{static} methods or compiler hints for deVirtualisation), since all method calls are resolved dynamically.

$K_{static}(\text{Ruby})=0$

\[
    O_{virt}(\text{Ruby}) = 0
\]

\textbf{Final metric value:}

\[
\begin{aligned}
    PDEI(\text{Ruby}) = 0.35 \cdot M_{scope}(\text{Ruby}) + 0.35 \cdot S_{disp}(\text{Ruby}) + 0.3 \cdot O_{virt}(\text{Ruby}) \\
    = 0.35 \cdot 0.2 + 0.35 \cdot 0 + 0.3 \cdot 0 \\
    = 0.07
\end{aligned}
\]

\subsubsection{State Model, Mutability, and Concurrency Guarantees}

\textbf{1) Default Mutability Score $M_{def}(L)$}

In Ruby, all values are objects, and most objects are \emph{mutable by default}, instances of classes such as \texttt{String}, \texttt{Array}, and \texttt{Hash} can be modified in place using methods like \texttt{<<}, \texttt{push}, or direct element assignment. Although Ruby provides mechanisms to enforce immutability, such immutability must be explicitly invoked and is not the default behaviour. Furthermore, there is no requirement for explicit annotation to enable mutability, as mutation is inherently permitted unless restricted.

\[
    M_{def}(\text{Ruby}) = 0
\]

\textbf{2) Static Reference Constraint Density $S_{dens}(L)$}

In Ruby, all variables are references to objects, and syntactic constructs include local variables, instance variables (\texttt{@x}), class variables (\texttt{@@x}), global variables (\texttt{\$x}), constants, and method parameters.

$N_{ref}(\text{Ruby}) = 6$ as distinct reference-related constructs explicitly represented in the language syntax. 

Ruby is a dynamically typed language with no static type system or compile-time enforcement of aliasing, mutability, or reference restrictions (all checks occur at runtime, and features like \texttt{freeze} are dynamic).

$N_{constr}(\text{Ruby}) = 0$ 

\[
    S_{dens}(\text{Ruby}) = 0
\]

\textbf{3) Concurrency Primitive Safety $P_{conc}(L)$}

First, Ruby does not provide compile-time guarantees of data race freedom, as it lacks a static type system capable of enforcing aliasing or ownership constraints. Moreover, while implementations such as MRI include a Global Interpreter Lock (GIL), this is an implementation detail rather than a language-level guarantee and does not eliminate race conditions in all cases (e.g., with native extensions or alternative runtimes). There are also no systematic runtime mechanisms that universally detect or prevent data races. Therefore, concurrency safety relies primarily on programmer discipline.

$S_{race}(\text{Ruby}) = 0$

Second, Ruby’s core concurrency primitives consist of low-level threads (\texttt{Thread}) and synchronisation mechanisms such as \texttt{Mutex}, without built-in high-level models like Actors or CSP channels in the core language (libraries may provide them, but they are not part of the language specification).

$S_{model}(\text{Ruby}) = 0$

\[
    P_{conc}(\text{Ruby})=0
\]

\textbf{Final metric value:}

\[
\begin{aligned}
    SCSI(\text{Ruby}) = 0.35 \cdot M_{def}(\text{Ruby}) + 0.25 \cdot S_{dens}(\text{Ruby}) + 0.4 \cdot P_{conc}(\text{Ruby}) \\
    = 0.35 \cdot 0 + 0.25 \cdot 0 + 0.4 \cdot 0 = 0
\end{aligned}
\]

\subsubsection{Support for Formal Specification and Verifiability}

\textbf{1) Native Contract Constructs Score $C_{nat}(L)$}

Ruby does not support built-in preconditions, postconditions, and invariants mechanisms for contract programming. The language does not provide special keywords and constructions for declaring and verifying. Users are able to implement checks during runtime using raise, methods, and metaprogramming

\[
    C_{nat}(\text{Ruby})=0
\]

\textbf{2) Verification Enforcement Level $V_{enf}(L)$}

Ruby does not support static verification. The language does not use static typing and compile-time correctness checks supports. Runtime checks are implemented using \texttt{raise}, assertion libraries, and reflection mechanisms. These tools allow developers to check the correctness of the program only during execution.

\[
    V_{enf}(\text{Ruby})=0.5
\]

\textbf{Final metric value:}

\[
\begin{aligned}
    FSVI(\text{Ruby}) = 0.45 \cdot C_{nat}(\text{Ruby}) + 0.55 \cdot V_{enf}(\text{Ruby}) \\
    = 0.45 \cdot 0 + 0.55 \cdot 0.5 \\
    = 0.275
\end{aligned}
\]
